\subsection{Kernel Selection by Stability and Causal Normalization}

The emergence of an exponentially screened memory kernel can be motivated from minimal physical requirements.

Let the memory kernel depend on the causal separation
\begin{equation}
\Delta t = t-t',
\qquad
\Delta t \geq 0 .
\end{equation}

Causality requires
\begin{equation}
K(\Delta t)=0
\qquad
{\rm for}
\qquad
\Delta t<0 .
\end{equation}

Finite memory requires the normalization condition
\begin{equation}
\int_0^\infty K(\Delta t)\,d\Delta t < \infty .
\end{equation}

Stability requires that the influence of remote past states decreases monotonically,
\begin{equation}
\frac{dK}{d\Delta t}<0 .
\end{equation}

The simplest first-order stable decay law satisfying these conditions is
\begin{equation}
\frac{dK}{d\Delta t}
=
-\lambda K,
\qquad
\lambda>0 .
\end{equation}

Its solution is
\begin{equation}
K(\Delta t)
=
K_0 e^{-\lambda \Delta t}.
\end{equation}

Including causal support explicitly gives
\begin{equation}
K(t,t')
=
K_0 e^{-\lambda(t-t')}
\Theta(t-t') .
\end{equation}

Thus, the exponential kernel is not an arbitrary choice. It is the minimal causal, normalizable and stable memory response.

When mapped to cosmological redshift or causal-depth variables, this kernel generates the leading SBC/F3 infrared correction
\begin{equation}
F_{\rm SBC}(z)
\simeq
-\epsilon e^{-k_0(1+z)} .
\end{equation}

This corresponds to the phenomenological SBC/F3 form
\begin{equation}
F_{\rm SBC}(z)
=
-\epsilon
\exp\left[-(k_0(1+z))^n\right]
\end{equation}
with
\begin{equation}
n\simeq 1 .
\end{equation}

Therefore, the observational preference for \(n\simeq1\) can be interpreted as the signature of the minimal stable causal-memory regime.
